\chapter[Integral Lipat dan Teorema Lapangan Vektor]{Integral Lipat, Perubahan Variabel, dan Teorema Lapangan Vektor}

Sebuah pelat pendingin berbentuk cakram berjari-jari satu meter. Tebalnya
seragam, tetapi kerapatannya menurun dari pusat: $e^{-(x^{2}+y^{2})}$ kilogram
per meter persegi. Massanya adalah jumlah kerapatan di seluruh cakram: pustaka numerik memberi
$1.9858653038$ kilogram setelah $11907$ evaluasi fungsi.

Rekan kita menulis versi lain dari perhitungan yang sama. Ia pindah ke
koordinat polar, menuliskan
\[
\int_{0}^{2\pi}\!\!\int_{0}^{1} e^{-r^{2}}\,dr\,d\theta ,
\]
dan mendapat $4.6924344183$ kilogram. Programnya berjalan tanpa peringatan.
Angka itu $2.36$ kali lebih besar daripada yang pertama, dan tidak satu pun
program memperingatkan. Kesalahannya satu huruf: faktor $r$ hilang, sehingga luas
tiap petak kecil dihitung terlalu besar.

Kerusakan sekelas itu muncul juga di sisi lain bab ini. Hitung kerja
$F(x,y)=(-y,x)/(x^{2}+y^{2})$ mengelilingi lingkaran satuan: hasilnya
$6.2831853072$, padahal curl medan itu nol di setiap titik tempat ia
terdefinisi. Kalau curl nol berarti kerja tertutup nol, salah satu hitungan
itu bohong. Yang bohong bukan hitungannya, melainkan syarat yang tidak
diperiksa: medan itu tidak terdefinisi di titik asal, dan titik asal ada di
dalam lingkaran.

Bab 13 membangun integral pada selang, dan Bab 12 memberi Jacobian sebagai
matriks turunan. Keduanya belum cukup: belum ada cara menjumlahkan pada daerah
berdimensi lebih dari satu, aturan tentang apa yang terjadi pada luas petak
ketika koordinat diganti, dan arti tanda pada batas. Bab ini mengisi
ketiganya, lalu menunjukkan bahwa Green, Gauss, dan Stokes sebenarnya satu
pernyataan.

\begin{untukanak}
Bayangkan penjahit dengan pola baju di atas kertas berpetak. Polanya persegi
dan rapi; kain yang akan dipotong melengkung mengikuti badan. Kalau pola persegi
itu ditempelkan begitu saja, bajunya sempit di satu sisi dan longgar di sisi
lain, walaupun jumlah petaknya tetap sama.

Penjahit berpengalaman melakukan hal lain: setiap petak kecil akan meregang
atau mengkerut dengan perbandingan tertentu saat dipindahkan ke kain. Petak di
dekat pinggang meregang sedikit; petak di dekat bahu mengkerut. Ia mencatat
perbandingan itu, mengalikan luas petak dengannya, baru menjumlahkan.

Pekerjaan yang sama muncul pada bentuk lengkung: banyak benda, luas, atau berat.
Kalau perbandingan regangan lupa dicatat, jumlahnya tetap bilangan rapi, dan
justru karena rapi ia sulit dicurigai. Di bab ini kita akan tahu mengapa
perbandingan itu harus ada, dan bagaimana menghitungnya untuk bentuk bulat.
\end{untukanak}

\section{Dari persegi panjang ke daerah umum}
\tingkat{dasar}

Kita mulai dari benda paling sederhana: sebuah persegi panjang, dan sebuah
bilangan yang hendak dijumlahkan di seluruh permukaannya. Caranya sama dengan
Bab 13, hanya potongannya sekarang berbentuk petak, bukan potongan garis.

\begin{definisi}[Integral lipat dua]\index{integral lipat dua}
Misalkan $R=[a,b]\times[c,d]$ sebuah persegi panjang tertutup. Sebuah
\emph{partisi} $\mathcal{P}$ dari $R$ adalah pembagian $[a,b]$ menjadi
selang-selang dan $[c,d]$ menjadi selang-selang; hasilnya adalah $mn$ persegi
panjang kecil $R_{ij}$. \emph{Mesh} partisi adalah diagonal terpanjang di
antara petak-petak itu, dan ditulis $\|\mathcal{P}\|$.

Untuk $f$ terbatas pada $R$, tulis $M_{ij}=\sup_{R_{ij}}f$ dan
$m_{ij}=\inf_{R_{ij}}f$, lalu
\[
U(\mathcal{P},f)=\sum_{i,j}M_{ij}\,\mathrm{luas}(R_{ij}),
\qquad
L(\mathcal{P},f)=\sum_{i,j}m_{ij}\,\mathrm{luas}(R_{ij}).
\]
Fungsi $f$ disebut \emph{terintegralkan} pada $R$ bila supremum semua $L$ sama
dengan infimum semua $U$; nilai bersama itu adalah $\iint_{R}f\,dA$.
\end{definisi}

Definisi ini tidak menyebut urutan pengintegralan. Bila $f$ kontinu, jumlah
atas dan jumlah bawah bertemu karena beda terbesar $f$ pada satu petak
mengecil bersama mesh: pada himpunan kompak, fungsi kontinu kontinu seragam
(Bab 11).

Pertanyaan berikutnya praktis. Menjumlahkan pada petak berarti menghitung
sesuatu berdimensi dua, sedangkan tangan kita hanya bisa menghitung integral satu
variabel. Jembatannya: memotong menurut satu arah lebih dulu, lalu menjumlahkan
hasil potongan itu.

\begin{teorema}[Fubini]\index{teorema Fubini}
Jika $f$ kontinu pada $R=[a,b]\times[c,d]$, maka
\[
\iint_{R} f\,dA=\int_{a}^{b}\!\!\int_{c}^{d} f(x,y)\,dy\,dx
=\int_{c}^{d}\!\!\int_{a}^{b} f(x,y)\,dx\,dy .
\]
\end{teorema}

\begin{proof}
Bukti penuh menuntut pencocokan jumlah atas-bawah dua dimensi dengan integral
satu dimensi secara seragam, dan ada di Spivak, bab \emph{Integration},
Teorema 3-10. Gagasannya begini.

Untuk $f\ge 0$, integral lipat dua adalah volume benda di bawah grafik $f$.
Iris benda itu dengan bidang $x$ tetap. Irisannya berupa daerah di bawah kurva
$y\mapsto f(x,y)$, sehingga luasnya
\[
A(x)=\int_{c}^{d}f(x,y)\,dy .
\]
Volume benda adalah integral dari luas irisan itu:
\[
\iint_{R}f\,dA=\int_{a}^{b}A(x)\,dx=\int_{a}^{b}\!\!\int_{c}^{d}f(x,y)\,dy\,dx .
\]
Karena $f$ kontinu pada himpunan kompak, $A$ kontinu. Untuk $f$ bertanda apa
pun, pisahkan $f=f^{+}-f^{-}$; arah sebaliknya diberikan oleh irisan $y$
tetap.
\end{proof}

Yang penting dari teorema ini bukan rumusnya, melainkan izinnya: urutan boleh
ditukar. Izin itu gugur pada fungsi yang tidak terintegralkan, dan itu dibahas
nanti.

Benda yang diiris dalam bukti itu tampak pada Gambar~\ref{fig:kolom-tipis}:
tiap kolom menyumbang $f(x,y)\,\Delta A$.

\begin{bacakeras}
--- Berapa isi gunungan pasir ini?
--- Jumlahkan kolom-kolom tipisnya. Berdiri di sisi utara, iris gunungan menjadi
jalur, ukur tiap jalur, lalu jumlahkan.
--- Berapa hasilnya?
--- Delapan belas ember.
--- Sekarang berdiri di sisi timur, iris dengan arah lain, ukur, jumlahkan.
--- Berapa hasilnya?
--- Delapan belas ember juga.
--- Jalur utara-selatan dan jalur timur-barat tidak satu pun sama. Mengapa
hasilnya sama?
--- Karena gunungannya satu. Kita tidak menjumlahkan dua gunungan; kita mengiris
benda yang sama dengan dua pisau berbeda arah.
--- Jadi arah iris boleh ditukar?
--- Selama bendanya terukur. Kalau gunungannya berlubang di tengah, dua arah itu
masih sama; yang gagal nanti lubangnya, bukan urutannya.

Ulangi sekali, dengan suara biasa: arah iris berganti, isi tetap.
\end{bacakeras}

\begin{figure}[htbp]
\centering
\begin{tikzpicture}[scale=1.55,font=\small]
  \fill[layergrey!10] (0,0) -- (6.1,0) -- (7.1,0.90) -- (1.0,0.90) -- cycle;
  \draw[layergrey,thick] (0,0) -- (6.1,0) -- (7.1,0.90) -- (1.0,0.90) -- cycle;
  \fill[layerteal!12] (0.91,0.82) rectangle (1.43,1.38);
  \draw[layerteal!70,line width=0.4pt] (0.91,0.82) rectangle (1.43,1.38);
  \fill[layerteal!30] (0.91,1.38) -- (1.08,1.47) -- (1.60,1.47) -- (1.43,1.38) -- cycle;
  \draw[layerteal!70,line width=0.4pt] (0.91,1.38) -- (1.08,1.47) -- (1.60,1.47);
  \fill[layerteal!12] (1.67,0.82) rectangle (2.19,1.60);
  \draw[layerteal!70,line width=0.4pt] (1.67,0.82) rectangle (2.19,1.60);
  \fill[layerteal!30] (1.67,1.60) -- (1.83,1.69) -- (2.35,1.69) -- (2.19,1.60) -- cycle;
  \draw[layerteal!70,line width=0.4pt] (1.67,1.60) -- (1.83,1.69) -- (2.35,1.69);
  \fill[layerteal!12] (2.42,0.82) rectangle (2.93,1.97);
  \draw[layerteal!70,line width=0.4pt] (2.42,0.82) rectangle (2.93,1.97);
  \fill[layerteal!30] (2.42,1.97) -- (2.58,2.06) -- (3.10,2.06) -- (2.93,1.97) -- cycle;
  \draw[layerteal!70,line width=0.4pt] (2.42,1.97) -- (2.58,2.06) -- (3.10,2.06);
  \fill[layerteal!12] (3.17,0.82) rectangle (3.68,2.30);
  \draw[layerteal!70,line width=0.4pt] (3.17,0.82) rectangle (3.68,2.30);
  \fill[layerteal!30] (3.17,2.30) -- (3.33,2.39) -- (3.85,2.39) -- (3.68,2.30) -- cycle;
  \draw[layerteal!70,line width=0.4pt] (3.17,2.30) -- (3.33,2.39) -- (3.85,2.39);
  \fill[layerteal!12] (3.92,0.82) rectangle (4.43,2.32);
  \draw[layerteal!70,line width=0.4pt] (3.92,0.82) rectangle (4.43,2.32);
  \fill[layerteal!30] (3.92,2.32) -- (4.08,2.41) -- (4.60,2.41) -- (4.43,2.32) -- cycle;
  \draw[layerteal!70,line width=0.4pt] (3.92,2.32) -- (4.08,2.41) -- (4.60,2.41);
  \fill[layerteal!12] (4.67,0.82) rectangle (5.18,2.00);
  \draw[layerteal!70,line width=0.4pt] (4.67,0.82) rectangle (5.18,2.00);
  \fill[layerteal!30] (4.67,2.00) -- (4.83,2.09) -- (5.35,2.09) -- (5.18,2.00) -- cycle;
  \draw[layerteal!70,line width=0.4pt] (4.67,2.00) -- (4.83,2.09) -- (5.35,2.09);
  \fill[layerteal!12] (5.42,0.82) rectangle (5.93,1.62);
  \draw[layerteal!70,line width=0.4pt] (5.42,0.82) rectangle (5.93,1.62);
  \fill[layerteal!30] (5.42,1.62) -- (5.58,1.71) -- (6.10,1.71) -- (5.93,1.62) -- cycle;
  \draw[layerteal!70,line width=0.4pt] (5.42,1.62) -- (5.58,1.71) -- (6.10,1.71);
  \fill[layerteal!12] (6.17,0.82) rectangle (6.68,1.39);
  \draw[layerteal!70,line width=0.4pt] (6.17,0.82) rectangle (6.68,1.39);
  \fill[layerteal!30] (6.17,1.39) -- (6.33,1.48) -- (6.85,1.48) -- (6.68,1.39) -- cycle;
  \draw[layerteal!70,line width=0.4pt] (6.17,1.39) -- (6.33,1.48) -- (6.85,1.48);
  \fill[layerteal!12] (0.59,0.50) rectangle (1.11,1.07);
  \draw[layerteal!70,line width=0.4pt] (0.59,0.50) rectangle (1.11,1.07);
  \fill[layerteal!30] (0.59,1.07) -- (0.76,1.16) -- (1.28,1.16) -- (1.11,1.07) -- cycle;
  \draw[layerteal!70,line width=0.4pt] (0.59,1.07) -- (0.76,1.16) -- (1.28,1.16);
  \fill[layerteal!12] (1.34,0.50) rectangle (1.86,1.33);
  \draw[layerteal!70,line width=0.4pt] (1.34,0.50) rectangle (1.86,1.33);
  \fill[layerteal!30] (1.34,1.33) -- (1.51,1.42) -- (2.03,1.42) -- (1.86,1.33) -- cycle;
  \draw[layerteal!70,line width=0.4pt] (1.34,1.33) -- (1.51,1.42) -- (2.03,1.42);
  \fill[layerteal!12] (2.09,0.50) rectangle (2.61,1.75);
  \draw[layerteal!70,line width=0.4pt] (2.09,0.50) rectangle (2.61,1.75);
  \fill[layerteal!30] (2.09,1.75) -- (2.26,1.84) -- (2.78,1.84) -- (2.61,1.75) -- cycle;
  \draw[layerteal!70,line width=0.4pt] (2.09,1.75) -- (2.26,1.84) -- (2.78,1.84);
  \fill[layerteal!12] (2.84,0.50) rectangle (3.36,2.14);
  \draw[layerteal!70,line width=0.4pt] (2.84,0.50) rectangle (3.36,2.14);
  \fill[layerteal!30] (2.84,2.14) -- (3.01,2.23) -- (3.53,2.23) -- (3.36,2.14) -- cycle;
  \draw[layerteal!70,line width=0.4pt] (2.84,2.14) -- (3.01,2.23) -- (3.53,2.23);
  \fill[layerteal!12] (3.59,0.50) rectangle (4.11,2.15);
  \draw[layerteal!70,line width=0.4pt] (3.59,0.50) rectangle (4.11,2.15);
  \fill[layerteal!30] (3.59,2.15) -- (3.76,2.24) -- (4.28,2.24) -- (4.11,2.15) -- cycle;
  \draw[layerteal!70,line width=0.4pt] (3.59,2.15) -- (3.76,2.24) -- (4.28,2.24);
  \fill[layerteal!12] (4.34,0.50) rectangle (4.86,1.79);
  \draw[layerteal!70,line width=0.4pt] (4.34,0.50) rectangle (4.86,1.79);
  \fill[layerteal!30] (4.34,1.79) -- (4.51,1.88) -- (5.03,1.88) -- (4.86,1.79) -- cycle;
  \draw[layerteal!70,line width=0.4pt] (4.34,1.79) -- (4.51,1.88) -- (5.03,1.88);
  \fill[layerteal!12] (5.09,0.50) rectangle (5.61,1.35);
  \draw[layerteal!70,line width=0.4pt] (5.09,0.50) rectangle (5.61,1.35);
  \fill[layerteal!30] (5.09,1.35) -- (5.26,1.44) -- (5.78,1.44) -- (5.61,1.35) -- cycle;
  \draw[layerteal!70,line width=0.4pt] (5.09,1.35) -- (5.26,1.44) -- (5.78,1.44);
  \fill[layerteal!12] (5.84,0.50) rectangle (6.36,1.08);
  \draw[layerteal!70,line width=0.4pt] (5.84,0.50) rectangle (6.36,1.08);
  \fill[layerteal!30] (5.84,1.08) -- (6.01,1.17) -- (6.53,1.17) -- (6.36,1.08) -- cycle;
  \draw[layerteal!70,line width=0.4pt] (5.84,1.08) -- (6.01,1.17) -- (6.53,1.17);
  \fill[layerteal!12] (0.26,0.17) rectangle (0.78,0.73);
  \draw[layerteal!70,line width=0.4pt] (0.26,0.17) rectangle (0.78,0.73);
  \fill[layerteal!30] (0.26,0.73) -- (0.43,0.82) -- (0.95,0.82) -- (0.78,0.73) -- cycle;
  \draw[layerteal!70,line width=0.4pt] (0.26,0.73) -- (0.43,0.82) -- (0.95,0.82);
  \fill[layerteal!12] (1.02,0.17) rectangle (1.54,0.95);
  \draw[layerteal!70,line width=0.4pt] (1.02,0.17) rectangle (1.54,0.95);
  \fill[layerteal!30] (1.02,0.95) -- (1.19,1.04) -- (1.71,1.04) -- (1.54,0.95) -- cycle;
  \draw[layerteal!70,line width=0.4pt] (1.02,0.95) -- (1.19,1.04) -- (1.71,1.04);
  \fill[layerteal!12] (1.76,0.17) rectangle (2.29,1.32);
  \draw[layerteal!70,line width=0.4pt] (1.76,0.17) rectangle (2.29,1.32);
  \fill[layerteal!30] (1.76,1.32) -- (1.93,1.41) -- (2.46,1.41) -- (2.29,1.32) -- cycle;
  \draw[layerteal!70,line width=0.4pt] (1.76,1.32) -- (1.93,1.41) -- (2.46,1.41);
  \fill[layerteal!12] (2.51,0.17) rectangle (3.04,1.65);
  \draw[layerteal!70,line width=0.4pt] (2.51,0.17) rectangle (3.04,1.65);
  \fill[layerteal!30] (2.51,1.65) -- (2.69,1.74) -- (3.21,1.74) -- (3.04,1.65) -- cycle;
  \draw[layerteal!70,line width=0.4pt] (2.51,1.65) -- (2.69,1.74) -- (3.21,1.74);
  \fill[layerteal!12] (3.26,0.17) rectangle (3.79,1.67);
  \draw[layerteal!70,line width=0.4pt] (3.26,0.17) rectangle (3.79,1.67);
  \fill[layerteal!30] (3.26,1.67) -- (3.44,1.76) -- (3.96,1.76) -- (3.79,1.67) -- cycle;
  \draw[layerteal!70,line width=0.4pt] (3.26,1.67) -- (3.44,1.76) -- (3.96,1.76);
  \fill[layerteal!12] (4.01,0.17) rectangle (4.53,1.35);
  \draw[layerteal!70,line width=0.4pt] (4.01,0.17) rectangle (4.53,1.35);
  \fill[layerteal!30] (4.01,1.35) -- (4.18,1.44) -- (4.70,1.44) -- (4.53,1.35) -- cycle;
  \draw[layerteal!70,line width=0.4pt] (4.01,1.35) -- (4.18,1.44) -- (4.70,1.44);
  \fill[layerteal!12] (4.76,0.17) rectangle (5.28,0.97);
  \draw[layerteal!70,line width=0.4pt] (4.76,0.17) rectangle (5.28,0.97);
  \fill[layerteal!30] (4.76,0.97) -- (4.93,1.06) -- (5.45,1.06) -- (5.28,0.97) -- cycle;
  \draw[layerteal!70,line width=0.4pt] (4.76,0.97) -- (4.93,1.06) -- (5.45,1.06);
  \fill[layerteal!12] (5.51,0.17) rectangle (6.03,0.74);
  \draw[layerteal!70,line width=0.4pt] (5.51,0.17) rectangle (6.03,0.74);
  \fill[layerteal!30] (5.51,0.74) -- (5.68,0.83) -- (6.20,0.83) -- (6.03,0.74) -- cycle;
  \draw[layerteal!70,line width=0.4pt] (5.51,0.74) -- (5.68,0.83) -- (6.20,0.83);
  \fill[accent!42] (2.69,0.17) rectangle (2.86,1.65);
  \draw[accent,line width=1pt] (2.69,0.17) rectangle (2.86,1.65);
  \fill[accent!55] (2.69,1.65) -- (2.78,1.70) -- (2.96,1.70) -- (2.86,1.65) -- cycle;
  \draw[accent,line width=1pt] (2.69,1.65) -- (2.78,1.70) -- (2.96,1.70);
  \draw[accent,densely dashed] (2.86,0.16) -- (3.42,-0.26);
  \node[accent,anchor=north west,align=left]
    at (3.50,-0.30) {satu kolom tipis: alas $\Delta A$,\\ tinggi $f(x,y)$, isi $\approx f(x,y)\,\Delta A$};
  \node[layergrey,anchor=north west] at (-0.05,-0.08) {alas $R$};
  \node[layergrey,anchor=north] at (5.95,-0.02) {$x$};
  \node[layergrey,anchor=west] at (7.15,0.55) {$y$};
\end{tikzpicture}
\caption{Benda di bawah grafik $f$ di atas persegi panjang $R$, dipotong menjadi
kolom-kolom tipis. Setiap kolom hampir berupa balok dengan alas $\Delta A$ dan tinggi
$f(x,y)$, sehingga isinya $f(x,y)\,\Delta A$; menjumlahkan seluruh kolom memberi
$\iint_R f\,dA$. Kolom yang ditandai merah adalah satu suku dalam jumlah itu.
Mengiris dari arah $x$ maupun dari arah $y$ menghasilkan kolom yang berbeda, tetapi
jumlah isinya sama.}
\label{fig:kolom-tipis}
\end{figure}

Daerah yang muncul di soal jarang berupa persegi panjang. Kita menanganinya
dengan satu langkah tambahan: perluas fungsinya dengan nol di luar daerah.

\begin{lemma}[Daerah di antara dua grafik]\index{daerah bertipe I}
Misalkan $g_{1},g_{2}\colon[a,b]\to\mathbb{R}$ kontinu dengan $g_{1}\le
g_{2}$, dan
\[
D=\{(x,y): a\le x\le b,\ g_{1}(x)\le y\le g_{2}(x)\}.
\]
Jika $f$ kontinu pada $D$, maka
\[
\iint_{D}f\,dA=\int_{a}^{b}\!\!\int_{g_{1}(x)}^{g_{2}(x)}f(x,y)\,dy\,dx .
\]
\end{lemma}

\begin{proof}
Ambil persegi panjang $R$ yang memuat $D$, dan perluas $f$ dengan nilai $0$ di
$R\setminus D$. Batas $D$ dapat ditutup oleh berhingga persegi panjang dengan
luas total kurang dari $\varepsilon$: pada grafik, kekontinuan seragam pada
selang kompak memberi partisi yang goyangan $g$ kecil di tiap potongan.
Himpunan yang dapat ditutup begitu tidak mengubah nilai integral, sehingga
\[
\iint_{D}f\,dA=\iint_{R} f\cdot\mathbf{1}_{D}\,dA .
\]
Menurut Fubini, integral kanan dihitung dengan mengintegralkan $x$ lebih dulu;
untuk $x$ tetap, integran $\mathbf{1}_{D}f$ bernilai $f$ hanya pada
$g_{1}(x)\le y\le g_{2}(x)$.
\end{proof}

Daerah bertipe I cukup untuk hampir semua soal di bab ini; batas yang memotong
garis tegak lebih dari dua kali dipecah lalu dijumlahkan. Cakram paling penting:
di Cartesius ia menuntut dua potongan, di polar cukup satu.

\begin{untukanak}
Sekarang ambil sebuah ruangan tertutup dengan beberapa tungku di dalamnya.
Setiap tungku mengeluarkan panas. Pertanyaannya: berapa panas yang keluar dari
ruangan itu setiap detik, dijumlahkan dari seluruh dinding, lantai, dan atap?

Cara pertama menghitungnya sulit: berkeliling ke setiap dinding, mengukur aliran
panas di tiap titik, lalu menjumlahkan. Cara kedua jauh lebih mudah: hitung
berapa panas yang dihasilkan semua tungku di dalam. Kalau tidak ada panas yang
hilang lewat celah lain, kedua angka harus sama: jumlah di dalam sama dengan
jumlah yang menembus kulit.

Kalimat itu terdengar terlalu bagus untuk benar. Kesamaan itu hanya berlaku
kalau kulit ruangan benar-benar tertutup dan arah "keluar" ditentukan dengan tegas. Kalau ada satu lubang
kecil yang tidak dihitung, atau kalau kita salah menandai mana sisi luar dan
mana sisi dalam, angkanya meleset walaupun semua penjumlahan kita benar.

Di bab ini kita akan tahu mengapa jumlah di dalam bisa dibaca dari kulitnya,
untuk bentuk lengkung apa pun, dan bagaimana satu lubang kecil dapat merusak
seluruh kesimpulan.
\end{untukanak}

\section{Perubahan variabel dan determinan Jacobian}
\tingkat{menengah}

Mengganti koordinat berarti mengganti petak: petak persegi di bidang $(u,v)$
menjadi petak melengkung di bidang $(x,y)$, dan luasnya berubah. Kita butuh
satu bilangan di tiap titik.

\begin{definisi}[Perubahan variabel]
\index{perubahan variabel}\index{Jacobian} Misalkan $D^{*}$ daerah terbuka di
$\mathbb{R}^{2}$ dan $T\colon D^{*}\to\mathbb{R}^{2}$,
$T(u,v)=(x(u,v),y(u,v))$, fungsi dengan turunan parsial kontinu. $T$ disebut
\emph{perubahan variabel} yang sah bila $T$ satu-satu pada bagian dalam
$D^{*}$ dan determinan Jacobiannya
\[
\det J_{T}(u,v)=\frac{\partial x}{\partial u}\frac{\partial y}{\partial v}
-\frac{\partial x}{\partial v}\frac{\partial y}{\partial u}
\]
tidak nol di sana. Matriks $J_{T}$ sendiri sudah kita kenal di Bab 12.
\end{definisi}

Kalimat sehari-harinya: ia angka yang memberi tahu berapa kali sebuah petak
kecil membesar ketika dipindahkan oleh $T$. Nilai $3$ berarti petak menjadi tiga kali lebih luas, $1/2$ berarti menyusut
setengah, dan nol berarti petak gepeng menjadi garis dan informasinya hilang.

Gambar~\ref{fig:juring-jacobian} memperlihatkan satu petak dan bayangannya:
luas $0.04$ menjadi $0.064$, dan angka $1.6$ itu adalah $|{\det J}|$ di tengah
petak.

\begin{figure}[htbp]
\centering
\begin{tikzpicture}[font=\footnotesize]
  \begin{scope}[x=1.9cm,y=1.9cm]
    \fill[layergrey!7] (1,0) rectangle (2.4,1.2);
    \draw[layergrey!40,very thin] (1,0) grid[xstep=0.2,ystep=0.2] (2.4,1.2);
    \draw[->,ink] (0.88,0) -- (2.62,0) node[right] {$u$};
    \draw[->,ink] (1,-0.08) -- (1,1.42) node[above] {$v$};
    \fill[layerindigo!30] (1.5,0.6) rectangle (1.7,0.8);
    \draw[layerindigo,thick] (1.5,0.6) rectangle (1.7,0.8);
    \node[layerindigo!85!black,anchor=north,align=center,fill=white,inner sep=1.2pt]
      at (1.60,0.55) {petak $0.2\times0.2$\\ luas $0.04$};
    \node[layergrey,anchor=north] at (1.7,-0.10) {bidang $(u,v)$};
  \end{scope}
  \draw[-{Latex[length=2.4mm]},ink,thick] (4.62,1.14) -- (5.85,1.14);
  \node[ink,above] at (5.23,1.18) {$T$};
  \node[ink,below,font=\scriptsize] at (5.23,1.06) {$\times\,1.6$};
  \begin{scope}[xshift=5.9cm,x=1.9cm,y=1.9cm]
    \foreach \r in {1,1.2,1.4,1.6,1.8,2,2.2,2.4}{
      \draw[layergrey!40,very thin] (\r,0) arc (0:69:\r);}
    \foreach \a in {0,11.5,23,34.5,46,57.5,69}{
      \draw[layergrey!40,very thin] (0,0) -- (\a:2.4);}
    \draw[->,ink] (-0.08,0) -- (2.62,0) node[right] {$x$};
    \draw[->,ink] (0,-0.08) -- (0,2.52) node[above] {$y$};
    \fill[accent!30] (34.4:1.5) arc (34.4:45.8:1.5) -- (45.8:1.7)
      arc (45.8:34.4:1.7) -- cycle;
    \draw[accent,thick] (34.4:1.5) arc (34.4:45.8:1.5) -- (45.8:1.7)
      arc (45.8:34.4:1.7) -- cycle;
    \node[accent,anchor=west,align=left] at (1.62,1.42)
      {juring melengkung\\ luas $0.064=1.6\times0.04$};
    \node[layergrey,anchor=north] at (1.6,-0.10) {bidang $(x,y)$};
  \end{scope}
\end{tikzpicture}
\caption{Satu petak kecil di bidang $(u,v)$ dan bayangannya di bidang $(x,y)$ di bawah
$T(u,v)=(u\cos v,\;u\sin v)$. Petak $0.2\times0.2$ berluas $0.04$ menjadi juring dengan
luas $0.064$, yaitu $1.6$ kali lipat, persis nilai $|{\det J}|=u$ di tengah petak
($u=1.6$). Perhatikan bahwa sisinya ikut melengkung: yang berubah bukan hanya ukuran,
melainkan juga bentuk, dan satu angka di tiap titik cukup untuk mencatat perubahannya.}
\label{fig:juring-jacobian}
\end{figure}

\begin{tangan}
Ambil kertas persegi $12$ cm kali $12$ cm, penggaris, pensil, gunting, dan
selotip. Gambar petak $1$ cm di seluruh kertas; ada $144$ petak. Potong kertas
menjadi dua belas jalur tegak, masing-masing $1$ cm kali $12$ cm. Pada setiap
jalur, tandai titik $3$ cm ke kanan di tepi atas, tarik garis dari sudut kiri
bawah ke titik itu, lalu potong sepanjang garis. Setiap jalur kini dua potongan
miring. Susun semua potongan menjadi satu jajar genjang yang rapat, lalu
rekatkan.

Lihat yang terjadi: bingkai luarnya $15$ cm kali $12$ cm, lebih lebar $3$ cm
daripada kertas asal, tetapi jajar genjang di dalamnya tetap berukuran alas
$12$ cm dan tinggi $12$ cm. Ukur keduanya dengan penggaris, tegak lurus, bukan
sepanjang sisi miring: hasilnya $144$ cm persegi, sama dengan luas kertas
sebelum dipotong. Kalau angkamu lain, yang salah hampir pasti tinggi.
\end{tangan}

\begin{teorema}[Perubahan variabel pada integral lipat dua]
\index{teorema perubahan variabel} Misalkan $T$ perubahan variabel yang sah
dari $D^{*}$ ke $D=T(D^{*})$, dan $f$ kontinu pada $D$. Maka
\[
\iint_{D}f(x,y)\,dx\,dy
=\iint_{D^{*}} f\bigl(T(u,v)\bigr)\,\bigl|\det J_{T}(u,v)\bigr|\,du\,dv .
\]
\end{teorema}

\begin{proof}
Kita buktikan gagasannya; bukti penuh ada di Spivak, bab \emph{Integration},
bagian perubahan variabel, dan di Rudin, bab \emph{Integration of Differential
Forms}.

Partisi $D^{*}$ menjadi petak kecil. Pada petak di sekitar $(u_{0},v_{0})$,
peta $T$ hampir linear: $T(u_{0}+h,v_{0}+k)\approx
T(u_{0},v_{0})+J_{T}(u_{0},v_{0})(h,k)$. Petak itu dipetakan menjadi jajar
genjang bersisi $h\,\partial T/\partial u$ dan $k\,\partial T/\partial v$,
yang luasnya dihitung dengan determinan (Bab 7):
\[
\mathrm{luas}
=\Bigl|h\,k\,\det\Bigl(\frac{\partial T}{\partial u},
\frac{\partial T}{\partial v}\Bigr)\Bigr|
=|h\,k|\,|\det J_{T}(u_{0},v_{0})| .
\]
Setiap petak menyumbang $f(T(u_{0},v_{0}))$ kali luas asalnya kali
$|\det J_{T}|$, dan jumlahnya adalah jumlah Riemann untuk integral di ruas
kanan; karena turunan $T$ kontinu pada himpunan kompak, galat aproksimasi linear
mengecil seragam bersama mesh.
\end{proof}

Nilai mutlak pada determinan bukan hiasan: determinan bisa negatif, luas tidak.
Tanda negatifnya menyimpan arah; pada integral lipat dua arah itu dibuang, tetapi
pada integral garis dan permukaan nanti tanda itu menentukan jawaban.

Sekarang determinan untuk dua sistem koordinat yang paling sering dipakai,
mulai dari polar.

\begin{proposisi}[Koordinat polar]\index{koordinat polar}
Dengan $x=r\cos\theta$, $y=r\sin\theta$, $r>0$, berlaku
\[
\det J=\begin{vmatrix}\cos\theta & -r\sin\theta\\ \sin\theta & r\cos\theta
\end{vmatrix}=r\cos^{2}\theta+r\sin^{2}\theta=r ,
\]
sehingga
$\iint_{D}f\,dx\,dy=\iint_{D^{*}}f(r\cos\theta,r\sin\theta)\,r\,dr\,d\theta$.
\end{proposisi}

Faktor $r$ itulah yang hilang pada program di pembukaan bab. Pada $r=1$
nilainya $1$, tetapi pada $r=0.1$ nilainya $0.1$: petak di dekat pusat
menyumbang sepuluh kali lebih kecil daripada dugaan yang mengabaikannya, dan
kesalahan totalnya bisa lebih dari dua kali.

\begin{proposisi}[Koordinat silinder dan bola]
\index{koordinat silinder}\index{koordinat bola} Untuk koordinat silinder
$x=r\cos\theta$, $y=r\sin\theta$, $z=z$, determinan Jacobiannya $r$. Untuk
koordinat bola
\[
x=\rho\sin\varphi\cos\theta,\quad
y=\rho\sin\varphi\sin\theta,\quad
z=\rho\cos\varphi ,
\]
dengan $\rho\ge0$, $0\le\varphi\le\pi$, $0\le\theta<2\pi$, determinan
Jacobiannya $\rho^{2}\sin\varphi$.
\end{proposisi}

\begin{proof}
Untuk silinder, kolom ketiga $(0,0,1)$ membuat determinan tinggal determinan
blok polar berukuran $2\times2$, yaitu $r$.

Untuk bola, hitung determinan $3\times3$ lewat ekspansi kofaktor. Dengan
$s_{\varphi}=\sin\varphi$, $c_{\varphi}=\cos\varphi$,
$s_{\theta}=\sin\theta$, $c_{\theta}=\cos\theta$, matriksnya
\[
J=\begin{pmatrix}
s_{\varphi}c_{\theta} & \rho c_{\varphi}c_{\theta} & -\rho s_{\varphi}s_{\theta}\\
s_{\varphi}s_{\theta} & \rho c_{\varphi}s_{\theta} & \rho s_{\varphi}c_{\theta}\\
c_{\varphi} & -\rho s_{\varphi} & 0
\end{pmatrix}.
\]
Entri baris ketiga kolom ketiga nol, jadi ekspansi kolom ketiga menyisakan dua
suku:
\[
\det J=\rho^{2}s_{\varphi}s_{\theta}^{2}+\rho^{2}s_{\varphi}c_{\theta}^{2}
=\rho^{2}s_{\varphi}\bigl(s_{\theta}^{2}+c_{\theta}^{2}\bigr)=\rho^{2}s_{\varphi}.
\]
\end{proof}

Faktor $\rho^{2}\sin\varphi$ nol di kutub $\varphi=0$ dan $\varphi=\pi$,
karena di sana seluruh bola $\rho$ runtuh menjadi satu titik dan peta berhenti
satu-satu; itu sebabnya kutub dikecualikan dari $D^{*}$. Ia paling besar di
ekuator, tempat perubahan kecil $\varphi$ memindahkan titik paling jauh.

\begin{resep}
\begin{enumerate}
\item Tulis $T$ eksplisit: $x=x(u,v)$, $y=y(u,v)$, dan periksa bahwa $T$
  satu-satu pada bagian dalam $D^{*}$.
\item Hitung empat turunan parsial, susun matriks $2\times2$, lalu hitung
  $\det J(u,v)$. Simpan tandanya; nilai mutlak dipasang nanti.
\item Terjemahkan batas: gambar $D$ di bidang $(x,y)$, tanyakan nilai $(u,v)$
  mana yang mendarat di sana, lalu tulis $D^{*}$.
\item Ganti seluruh $x,y$ pada $f$ dengan $x(u,v),y(u,v)$, kalikan
  $|\det J(u,v)|$, lalu tutup dengan $du\,dv$ sesuai urutan batas.
\end{enumerate}
Peringatan urutan: langkah 2 dan 3 tidak boleh ditukar. Determinan harus sudah
berbentuk $(u,v)$ sebelum batas dipasang. Kalau batas ditulis lebih dulu, faktor itu tertinggal dalam $x$ dan $y$ dan
integralnya tidak dapat dihitung.
\end{resep}

\kamus{determinan Jacobian}{Jacobian determinant}{Angka yang mengatakan berapa kali petak kecil meregang ketika koordinat diganti.}

\kamus{integral lipat dua}{double integral}{Menjumlahkan nilai pada petak kecil di sebuah permukaan, lalu mempersempitnya.}

\kamus{orientasi}{orientation}{Kesepakatan tentang arah putar mana yang dihitung positif.}

\section{Integral garis dan medan konservatif}
\tingkat{menengah}

Sekarang kita berhenti menjumlahkan pada permukaan dan mulai menjumlahkan
sepanjang jalan. Bayangkan sebuah benda kecil diseret sepanjang lintasan
melengkung oleh gaya yang arah dan besarnya berubah di setiap titik. Kerja
gaya itu adalah jumlah komponen gaya searah jalan dikali potongan panjang
jalan.

\begin{definisi}[Integral garis]\index{integral garis}
Misalkan $\gamma\colon[a,b]\to\mathbb{R}^{2}$ kurva dengan turunan kontinu
sepotong-potong, dan $F=(P,Q)$ medan vektor kontinu pada lintasannya.
\emph{Integral garis} $F$ sepanjang $\gamma$ adalah
\[
\int_{\gamma}F\cdot d\mathbf{r}
=\int_{a}^{b}\Bigl(P(\gamma(t))\,\gamma_{1}'(t)
+Q(\gamma(t))\,\gamma_{2}'(t)\Bigr)\,dt .
\]
Bila $\gamma$ tertutup, artinya $\gamma(a)=\gamma(b)$, integral itu disebut
\emph{sirkulasi} $F$ mengelilingi $\gamma$.
\end{definisi}

Rumus itu tidak memilih titik awal, hanya arah. Mengganti parameter dengan
laju berbeda tidak mengubah nilainya, karena substitusi (Bab 13) menyerap
faktor laju; membalik arah mengubah tanda. Jadi integral garis berorientasi,
dan itu perbedaan pertamanya dengan integral lipat dua.

Gambar~\ref{fig:sirkulasi} menunjukkan medan berputar seperti itu, dengan satu
kurva tertutup untuk mengukur berapa banyak putaran yang tertangkap.

\begin{figure}[htbp]
\centering
\begin{tikzpicture}[scale=1.45,font=\footnotesize]
  \draw[->,layergrey] (-2.85,0) -- (2.85,0) node[right] {$x$};
  \draw[->,layergrey] (0,-2.85) -- (0,2.85) node[above] {$y$};
  \foreach \x in {-2.4,-1.6,-0.8,0,0.8,1.6,2.4}{
    \foreach \y in {-2.4,-1.6,-0.8,0,0.8,1.6,2.4}{
      \draw[-{Latex[length=1.2mm]},layerteal!75]
        (\x,\y) -- ({\x-0.30*\y},{\y+0.30*\x});}}
  \draw[accent,very thick] (0,0) circle (1.75);
  \draw[-{Latex[length=2.8mm]},accent,very thick]
    (1.126,-1.341) -- (1.340,-1.126);
  \node[accent,anchor=north west,align=left] at (0.28,-1.86)
    {kurva tertutup, arah positif};
  \node[ink,align=center,fill=white,fill opacity=0.85,text opacity=1,inner sep=2pt] at (0,0.95)
    {$\displaystyle\oint_{\gamma}F\cdot d\mathbf{r}=2\pi$};
  \node[ink,font=\scriptsize,fill=white,fill opacity=0.85,text opacity=1,inner sep=1.5pt] at (0,0.42)
    {sedangkan $\operatorname{curl}F=2$ di tiap titik};
\end{tikzpicture}
\caption{Medan $F(x,y)=(-y,x)$ dengan panah di setiap titik kisi. Panahnya berputar
mengelilingi titik asal, dan besarnya bertambah sejauh titik itu dari pusat. Sirkulasi
mengelilingi lingkaran berjari-jari $1$ adalah $2\pi$, sama dengan luas $\pi$ dikali
$\operatorname{curl}F=2$: inilah Green untuk medan ini. Arah panah pada kurva adalah arah
positif, berlawanan arah jarum jam.}
\label{fig:sirkulasi}
\end{figure}

Kalau medan gaya berasal dari sebuah fungsi ketinggian, kerja hanya bergantung
pada titik awal dan titik akhir. Ini versi kurva dari teorema dasar kalkulus,
dan namanya mengikuti hal itu.

\begin{teorema}[Teorema dasar untuk integral garis]
\index{medan konservatif} Misalkan $f\colon U\to\mathbb{R}$ mempunyai turunan
parsial kontinu pada himpunan terbuka $U\subseteq\mathbb{R}^{2}$, dan $\gamma$
kurva dari $A$ ke $B$ di dalam $U$. Maka
\[
\int_{\gamma}\nabla f\cdot d\mathbf{r}=f(B)-f(A).
\]
\end{teorema}

\begin{proof}
Tulis $g(t)=f(\gamma(t))$. Menurut aturan rantai multivariabel (Bab 12),
\[
g'(t)=\nabla f(\gamma(t))\cdot\gamma'(t),
\]
asalkan $\gamma$ terdiferensialkan di $t$. Jadi integran pada integral garis
adalah turunan $g$, dan teorema dasar kalkulus (Bab 13) memberi
$\int_{a}^{b}g'(t)\,dt=g(b)-g(a)=f(B)-f(A)$. Untuk kurva yang hanya
terdiferensialkan sepotong-potong, jumlahkan pada tiap potongan; suku di titik
sambung saling menghapus.
\end{proof}

Medan berbentuk $\nabla f$ disebut \emph{konservatif}, dan $f$ potensialnya.
Akibat langsung teoremanya: pada medan konservatif, sirkulasi mengelilingi
lintasan tertutup selalu nol.

Kebalikannya juga benar.

\begin{teorema}[Kebalikan]
Misalkan $F$ kontinu pada himpunan terbuka terhubung
$U\subseteq\mathbb{R}^{2}$, dan integral garis $F$ tidak bergantung pada
lintasan di dalam $U$. Maka $F=\nabla f$ untuk suatu $f$, yaitu
\[
f(\mathbf{x})=\int_{\mathbf{x}_{0}}^{\mathbf{x}}F\cdot d\mathbf{r},
\]
dengan $\mathbf{x}_{0}\in U$ tetap.
\end{teorema}

\begin{proof}
Ambil $\mathbf{x}\in U$ dan $h$ kecil sehingga segmen dari $\mathbf{x}$ ke
$\mathbf{x}+h\mathbf{e}_{1}$ tinggal di dalam $U$. Karena lintasan tidak
penting, kita boleh menempuh jalan dari $\mathbf{x}_{0}$ ke $\mathbf{x}$ lalu
menambah segmen pendek itu, sehingga
\[
f(\mathbf{x}+h\mathbf{e}_{1})-f(\mathbf{x})
=\int_{0}^{1}F_{1}(\mathbf{x}+th\mathbf{e}_{1})\,h\,dt .
\]
Bagi dengan $h$ dan pakai kekontinuan: nilai rata-rata integran menuju
$F_{1}(\mathbf{x})$ ketika $h\to0$, jadi $\partial f/\partial x_{1}=F_{1}$,
dan argumen sama memberi $\partial f/\partial x_{2}=F_{2}$. Kedua turunan
parsial itu kontinu, maka $F=\nabla f$.
\end{proof}

Kalau diberi $F$, bagaimana tahu ia konservatif tanpa mencari potensial? Di
bidang ada uji cepat: $F=(P,Q)$ berpotensial hanya jika $\partial P/\partial
y=\partial Q/\partial x$, sebab kedua ruas adalah turunan kedua $f$ dan Clairaut
(Bab 12) menyamakannya. Arah sebaliknya menuntut domain tanpa lubang dan
dibuktikan lewat Green nanti. Medan pembukaan bab tandingannya: curl-nya nol,
tetapi ia bukan gradien apa pun.

\begin{contoh}[Dua cara menghitung kerja]
\index{potensial} Ambil $F(x,y)=(2xy,\ x^{2}+3y^{2})$ dan lintasan
$\gamma(t)=(t,2t^{2})$ untuk $0\le t\le1$, dari $(0,0)$ ke $(1,2)$.

Cara pertama, langsung. Di sini $dx=dt$ dan $dy=4t\,dt$, sehingga
\[
\int_{\gamma}F\cdot d\mathbf{r}
=\int_{0}^{1}\Bigl(2t\cdot 2t^{2}+(t^{2}+12t^{4})\cdot 4t\Bigr)dt
=\int_{0}^{1}\bigl(8t^{3}+48t^{5}\bigr)dt=2+8=10 .
\]
Cara kedua, lewat potensial. Periksa dulu bahwa $\partial P/\partial
y=2x=\partial Q/\partial x$, lalu tebak $f(x,y)=x^{2}y+y^{3}$. Turunannya
memang $\partial f/\partial x=2xy$ dan $\partial f/\partial y=x^{2}+3y^{2}$.
Karena itu
\[
\int_{\gamma}F\cdot d\mathbf{r}=f(1,2)-f(0,0)=(2+8)-0=10 .
\]
Kedua cara memberi $10$; kuadratur adaptif berselisih $1.78\times10^{-15}$.
\end{contoh}

\section{Teorema Green}
\tingkat{menengah}

Gagasan di kotak tentang ruangan dan tungku sekarang menjadi teorema: apa yang
terjadi di dalam sebuah daerah dapat dibaca dari tepinya.

Sebelum itu, satu kesepakatan arah.

\begin{definisi}[Orientasi positif]
\index{orientasi positif} Kurva tertutup sederhana $\gamma$ berorientasi
\emph{positif} bila daerah yang dibatasinya selalu berada di sebelah kiri
ketika kita menelusurinya. Untuk lingkaran, arah positif adalah berlawanan
arah jarum jam.
\end{definisi}

\begin{teorema}[Green]\index{teorema Green}
Misalkan $D$ daerah tertutup terbatas di $\mathbb{R}^{2}$ yang batasnya
$\partial D$ terdiri dari berhingga kurva tertutup sederhana yang mulus
sepotong-potong dan berorientasi positif. Misalkan $P,Q$ mempunyai turunan
parsial kontinu pada himpunan terbuka yang memuat $D$. Maka
\[
\oint_{\partial D}P\,dx+Q\,dy=\iint_{D}\Bigl(\frac{\partial Q}{\partial x}
-\frac{\partial P}{\partial y}\Bigr)dA .
\]
\end{teorema}

\begin{proof}
Cukup dibuktikan untuk daerah bertipe I sekaligus bertipe II, karena daerah
umum dipecah menjadi berhingga potongan semacam itu dan integralnya
dijumlahkan; sisi potongan yang bersebelahan menempuh arah berlawanan sehingga
saling menghapus.

Andaikan $D$ bertipe I dengan $g_{1},g_{2}$ mulus. Pada sisi tegak $dx=0$,
jadi sumbangannya nol dan
\[
\oint_{\partial D}P\,dx=\int_{a}^{b}P(x,g_{1}(x))\,dx+\int_{b}^{a}P(x,g_{2}(x))\,dx
=\int_{a}^{b}\bigl[P(x,g_{1})-P(x,g_{2})\bigr]dx .
\]
Menurut teorema dasar kalkulus,
$P(x,g_{1})-P(x,g_{2})=-\int_{g_{1}(x)}^{g_{2}(x)}\partial P/\partial y\,dy$,
sehingga
\[
\oint_{\partial D}P\,dx=-\int_{a}^{b}\!\!\int_{g_{1}(x)}^{g_{2}(x)}
\frac{\partial P}{\partial y}\,dy\,dx=-\iint_{D}\frac{\partial P}{\partial y}\,dA .
\]
Dengan cara yang sama pada $D$ bertipe II, tempat $dy=0$ pada sisi datar,
$\oint_{\partial D}Q\,dy=\iint_{D}\partial Q/\partial x\,dA$.
\end{proof}

\begin{contoh}[Green pada persegi satuan]
Ambil $P=x^{2}-y$, $Q=x+2y$, dan $D=[0,1]\times[0,1]$. Di sini $\partial
Q/\partial x-\partial P/\partial y=1-(-1)=2$, sehingga ruas kanan Green adalah
$2$ kali luas persegi, yaitu $2$.

Ruas kiri kita hitung pada empat sisi, berlawanan arah jarum jam:
\begin{align*}
\text{bawah }(y=0,\,dy=0):&\quad \int_{0}^{1}x^{2}\,dx=\tfrac13 ,\\
\text{kanan }(x=1,\,dx=0):&\quad \int_{0}^{1}(1+2y)\,dy=2 ,\\
\text{atas }(y=1,\,dy=0):&\quad \int_{1}^{0}(x^{2}-1)\,dx=\tfrac23 ,\\
\text{kiri }(x=0,\,dx=0):&\quad \int_{1}^{0}2y\,dy=-1 .
\end{align*}
Jumlahnya $\tfrac13+2+\tfrac23-1=2$, sama dengan ruas kanan. Kuadratur adaptif
untuk keempat sisi memberi $2.000000000000000$, dan integral lipat duanya juga
$2.000000000000000$.
\end{contoh}

Sekarang kesalahan tanda yang diminta. Ulangi hitungan yang sama searah jarum
jam. Setiap sisi ditempuh terbalik, jadi setiap suku berganti tanda dan
jumlahnya menjadi $-2$; ruas kanan Green tidak berubah karena daerahnya sama.
Jawaban yang seharusnya $2$ menjadi $-2$: bukan meleset sedikit, melainkan salah
tanda sepenuhnya, dengan selisih $4$ yang tidak dapat disalahkan pada
pembulatan.

Jadi orientasi bukan kesepakatan sopan: rumus Green hanya benar untuk satu
arah putar, dan arah itu harus diperiksa sebelum menghitung.

Dengan Green di tangan, arah sebaliknya dari uji curl di bagian sebelumnya
dapat dibuktikan.

\begin{akibat}[Uji curl di domain tanpa lubang]\index{uji curl}
Misalkan $U\subseteq\mathbb{R}^{2}$ terbuka dan \emph{terhubung sederhana},
artinya setiap kurva tertutup di $U$ dapat disusutkan ke satu titik di dalam
$U$. Misalkan $P,Q$ mempunyai turunan parsial kontinu pada $U$ dan $\partial
P/\partial y=\partial Q/\partial x$ di seluruh $U$. Maka $F=(P,Q)$ konservatif
pada $U$.
\end{akibat}

\begin{proof}
Ambil kurva tertutup sederhana $\gamma$ di $U$ dan daerah $D$ yang
dibatasinya. Karena $U$ terhubung sederhana, $D$ tinggal di dalam $U$, jadi
Green berlaku dan
\begin{align*}
\oint_{\gamma}P\,dx+Q\,dy
&=\iint_{D}\Bigl(\frac{\partial Q}{\partial x}-\frac{\partial P}{\partial y}\Bigr)dA\\
&=0 .
\end{align*}
Kurva tertutup yang memotong dirinya sendiri dipecah menjadi lengkung-lengkung
sederhana, dan tiap potongan memberi nol. Semua sirkulasi nol, maka integral
garis tidak bergantung lintasan, dan teorema kebalikan tadi memberi potensial.
\end{proof}

Syarat itu tidak bisa dibuang. Pada daerah berlubang, Green menuntut batas
dalam ikut dihitung; di situlah medan pembukaan bab lolos dari uji curl.

\section{Teorema Gauss dan Stokes}
\tingkat{lanjut}

Green berbicara tentang daerah di bidang dan tepinya. Di ruang berdimensi tiga
ada dua versi, keduanya memakai divergensi dan curl dari Bab 12.

\begin{definisi}[Divergensi dan curl]
\index{divergensi}\index{curl} Untuk medan $F=(F_{1},F_{2},F_{3})$ dengan
turunan parsial kontinu,
\[
\operatorname{div}F=\frac{\partial F_{1}}{\partial x}
+\frac{\partial F_{2}}{\partial y}+\frac{\partial F_{3}}{\partial z},
\qquad
\operatorname{curl}F=\Bigl(\frac{\partial F_{3}}{\partial y}
-\frac{\partial F_{2}}{\partial z},\
\frac{\partial F_{1}}{\partial z}-\frac{\partial F_{3}}{\partial x},\
\frac{\partial F_{2}}{\partial x}-\frac{\partial F_{1}}{\partial y}\Bigr).
\]
\end{definisi}

Divergensi mengukur berapa banyak medan lahir di sebuah titik: positif berarti
sumber, negatif penyerap. Curl mengukur putaran medan di sekitar titik itu.

\begin{teorema}[Gauss, atau teorema divergensi]\index{teorema Gauss}
Misalkan $E\subseteq\mathbb{R}^{3}$ daerah tertutup terbatas yang batasnya
$\partial E$ berupa permukaan mulus sepotong-potong, dan $\mathbf{n}$ vektor
normal satuan yang menunjuk keluar. Jika $F$ mempunyai turunan parsial kontinu
pada himpunan terbuka yang memuat $E$, maka
\[
\iiint_{E}\operatorname{div}F\,dV=\iint_{\partial E}F\cdot\mathbf{n}\,dA .
\]
\end{teorema}

\begin{proof}
Kita buktikan untuk kotak dulu.

Ambil $E=[a,b]\times[c,d]\times[e,f]$. Untuk suku ketiga,
\[
\iiint_{E}\frac{\partial F_{3}}{\partial z}\,dV
=\iint_{[a,b]\times[c,d]}\bigl[F_{3}(x,y,f)-F_{3}(x,y,e)\bigr]dx\,dy ,
\]
menurut Fubini dan teorema dasar kalkulus pada arah $z$. Ruas kanan itu tepat
fluks lewat sisi atas (normal $+\mathbf{k}$) dan sisi bawah (normal
$-\mathbf{k}$), karena di sana $\mathbf{n}\,dA=\pm\mathbf{k}\,dx\,dy$. Suku
lain memberi fluks lewat dua pasang sisi lainnya.

Untuk daerah umum, potong menjadi kotak-kotak kecil. Sisi dalam muncul dua kali
dengan normal berlawanan dan saling menghapus, sehingga yang tersisa hanya sisi
yang menyentuh batas luar; limit mesh menuju nol memberi pernyataan untuk batas
lengkung. Bukti penuh ada di Spivak, Teorema 5-8.
\end{proof}

\begin{contoh}[Gauss pada bola satuan]
Ambil $F(x,y,z)=(x^{3},y^{3},z^{3})$ dan $E$ bola satuan. Di sini
$\operatorname{div}F=3(x^{2}+y^{2}+z^{2})=3\rho^{2}$, sehingga dengan
determinan Jacobian bola
\[
\iiint_{E}3\rho^{2}\,dV=\int_{0}^{2\pi}\!\!\int_{0}^{\pi}\!\!\int_{0}^{1}
3\rho^{4}\sin\varphi\,d\rho\,d\varphi\,d\theta
=3\cdot2\pi\cdot2\cdot\frac15=\frac{12\pi}{5}.
\]
Ruas permukaan: pada bola satuan $\mathbf{n}=(x,y,z)$, jadi
$F\cdot\mathbf{n}=x^{4}+y^{4}+z^{4}$ dan
\[
\iint_{\partial E}(x^{4}+y^{4}+z^{4})\,dA=3\iint_{\partial E}z^{4}\,dA
=3\int_{0}^{2\pi}\!\!\int_{0}^{\pi}\cos^{4}\varphi\sin\varphi\,d\varphi\,d\theta
=3\cdot2\pi\cdot\frac25=\frac{12\pi}{5},
\]
dengan faktor $3$ dari simetri. Nilainya $7.5398223686155$, dan kuadratur pada
permukaan memberi angka yang sama sampai digit terakhir.
\end{contoh}

Teorema kedua berbicara tentang permukaan tidak tertutup yang batasnya kurva.
Rumusnya memerlukan aturan tangan kanan.

\begin{teorema}[Stokes]\index{teorema Stokes}
Misalkan $S$ permukaan mulus sepotong-potong berorientasi dengan normal
$\mathbf{n}$, dan $\partial S$ kurva batasnya dengan orientasi yang diberikan
aturan tangan kanan. Jika $F$ mempunyai turunan parsial kontinu pada himpunan
terbuka yang memuat $S$, maka
\[
\iint_{S}(\operatorname{curl}F)\cdot\mathbf{n}\,dA
=\oint_{\partial S}F\cdot d\mathbf{r}.
\]
\end{teorema}

\begin{proof}
Parametrisasi $S$ oleh $\mathbf{r}(u,v)$ atas daerah $D$ di bidang $(u,v)$,
dengan orientasi yang sesuai. Aturan rantai mengubah integrannya menjadi
turunan komponen $F$ terhadap $u$ dan $v$, dan setelah disusun ulang integral
permukaan itu menjadi
\[
\iint_{D}\Bigl(\frac{\partial \tilde{Q}}{\partial u}
-\frac{\partial \tilde{P}}{\partial v}\Bigr)du\,dv ,
\]
untuk dua fungsi $\tilde{P},\tilde{Q}$ yang dibentuk dari $F$ dan
$\mathbf{r}$. Itu tepat ruas kanan Green, jadi Green memberi integral garis
pada $\partial D$, yang tidak lain integral garis $F$ pada $\partial S$.
Aljabar penuhnya di Spivak, Teorema 5-9.
\end{proof}

Kalau $S$ datar di bidang $xy$, Stokes menjadi Green; kalau $S$ hanya sebuah
selang, ia menjadi teorema dasar kalkulus.

\begin{contoh}[Stokes pada belahan bola]
Ambil $F=(z,x,y)$, yang curl-nya $(1,1,1)$, dan $S$ belahan atas bola satuan
dengan normal keluar, sehingga batasnya lingkaran satuan di bidang $xy$
berlawanan jarum jam.

Ruas permukaan:
\[
\iint_{S}(1,1,1)\cdot\mathbf{n}\,dA=\iint_{S}(x+y+z)\,dA=\iint_{S}z\,dA
=\int_{0}^{2\pi}\!\!\int_{0}^{\pi/2}\cos\varphi\sin\varphi\,d\varphi\,d\theta
=2\pi\cdot\frac12=\pi ,
\]
karena suku $x$ dan $y$ saling menghapus oleh simetri belahan. Ruas batas, dengan
$\mathbf{r}(t)=(\cos t,\sin t,0)$:
\[
\oint_{\partial S}F\cdot d\mathbf{r}=\int_{0}^{2\pi}\cos^{2}t\,dt=\pi .
\]
Keduanya $3.1415926535898$; kuadratur memberi angka yang sama pada kedua ruas.
\end{contoh}

Sekarang kesalahan tanda pada Stokes. Balik normal permukaan menjadi ke dalam:
ruas permukaan menjadi $-\pi$, ruas batas tidak berubah selama arah putarnya
tetap, dan kesamaan gagal dengan selisih $2\pi$. Aturan tangan kanan yang
mengikat kedua ruas; membalik keduanya bersama memulihkan kesamaan.

\section{Bentuk diferensial: satu kalimat untuk semua}
\tingkat{lanjut}

Kita sudah melihat tiga teorema dengan tiga wajah. Ada bahasa yang
menuliskannya sebagai satu pernyataan, dan bahasanya bukan sekadar kerapian: ia
menjelaskan mengapa tanda orientasi muncul di mana-mana.

\begin{definisi}[Bentuk diferensial]
\index{bentuk diferensial} Sebuah \emph{bentuk-$k$} pada himpunan terbuka
$U\subseteq\mathbb{R}^{n}$ adalah jumlah suku berbentuk
$a(x)\,dx_{i_{1}}\wedge\cdots\wedge dx_{i_{k}}$, dengan $a$ fungsi mulus.
Perkalian baji $\wedge$ bersifat antisimetris:
\[
dx\wedge dy=-dy\wedge dx,\qquad dx\wedge dx=0 .
\]
Turunan luar $d$ bekerja pada koefisien dan menambah satu baji:
\[
d(a\,dx_{I})=da\wedge dx_{I},\qquad da=\sum_{j}\frac{\partial a}{\partial x_{j}}dx_{j}.
\]
\end{definisi}

Kalimat sehari-harinya: bentuk diferensial adalah benda yang siap
diintegralkan pada permukaan berdimensi tertentu, dan $d$ menaikkan dimensi satu
tingkat. Antisimetri membuat membalik orientasi otomatis membalik tanda;
kesulitan tanda sebelumnya tertanam di satu aturan itu.

\begin{teorema}[Stokes umum]
Misalkan $M$ manifold berorientasi berdimensi $k$ dengan batas, dan $\omega$
bentuk-$(k-1)$ dengan koefisien mulus. Maka
\[
\int_{\partial M}\omega=\int_{M}d\omega ,
\]
dengan orientasi $\partial M$ yang diinduksi oleh orientasi $M$.
\end{teorema}

Bukti lengkapnya dilewati: Spivak, \emph{Integration on Chains}, Teorema 4-13
untuk rantai, dan Teorema 5-5 pada \emph{Integration of Manifolds} untuk
manifold; juga Rudin, bab \emph{Integration of Differential Forms}. Di sana
Green, Gauss, dan Stokes klasik diturunkan sebagai kasus khusus.

Terjemahannya di ruang berdimensi tiga:

\begin{center}
\begin{tabular}{lll}
\hline
Bahasa bentuk & Medan vektor & Teorema \\
\hline
bentuk-0 $f$ & fungsi skalar & $df$ memberi gradien \\
bentuk-1 $P\,dx+Q\,dy+R\,dz$ & $F=(P,Q,R)$ & $d\omega$ memberi $\operatorname{curl}F$ \\
bentuk-2 $P\,dy\wedge dz+\cdots$ & $F=(P,Q,R)$ & $d\eta$ memberi $\operatorname{div}F$ \\
$\int_{\partial M}\omega=\int_{M}d\omega$ & & Green, Gauss, Stokes, FTC \\
\hline
\end{tabular}
\end{center}

Satu akibat kecil: karena $d(d\omega)=0$, turunan luar dua kali selalu nol.
Diterjemahkan, $\operatorname{curl}\nabla f=0$ dan
$\operatorname{div}\operatorname{curl}F=0$. Pemeriksaan simbolik untuk fungsi
dan medan umum memberi nol pada keenam komponen, jadi dua identitas yang tampak
kebetulan itu satu pernyataan.

\section{Kapan teorema ini gagal}
\tingkat{lanjut}

Semua teorema di bab ini kesamaan bersyarat. Bagian ini mengumpulkan empat
cara kesamaan itu rusak, dengan angka.

Kegagalan pertama adalah singularitas. Ambil $F(x,y)=(-y,x)/(x^{2}+y^{2})$
pada bidang tanpa titik asal. Hitung curl-nya:
\[
\frac{\partial Q}{\partial x}-\frac{\partial P}{\partial y}
=\frac{(x^{2}+y^{2})-2x^{2}}{(x^{2}+y^{2})^{2}}
+\frac{(x^{2}+y^{2})-2y^{2}}{(x^{2}+y^{2})^{2}}=0 .
\]
Beda hingga $h=10^{-5}$ di $(1,0)$ memberi $-2.01\times10^{-10}$, nol sampai
galat pemotongan. Tetapi sirkulasi
mengelilingi lingkaran satuan $2\pi=6.2831853072$, karena di sana
$F\cdot\gamma'=1$ di setiap titik. Jadi curl nol tidak berarti sirkulasi nol.

Yang dilanggar adalah syarat Green: $P$ dan $Q$ harus punya turunan parsial
kontinu pada himpunan terbuka yang memuat $D$, sedangkan keduanya tidak
terdefinisi di titik asal di dalam $D$, dan $D$ tidak terhubung sederhana. Pada cincin $1\le r\le2$ Green tetap berlaku dengan dua batas
berlawanan arah, dan integralnya $2\pi-2\pi=0$; nol itu benar karena di sana
medan berpotensial $\theta$, sudut polar.

Kegagalan kedua adalah pertukaran urutan tanpa keterintegralan. Ambil
\[
f(x,y)=\frac{x^{2}-y^{2}}{(x^{2}+y^{2})^{2}}
\]
pada persegi $(0,1)\times(0,1)$. Untuk $y>0$ tetap,
\[
\int_{0}^{1}f(x,y)\,dx=\Bigl[-\frac{x}{x^{2}+y^{2}}\Bigr]_{x=0}^{x=1}
=-\frac{1}{1+y^{2}},
\]
sehingga integral luar memberi $-\pi/4=-0.7853981634$; dengan urutan terbalik
hasilnya $+\pi/4$. Kedua integral berulang ada tetapi berbeda, jadi menurut
Fubini integral lipat duanya tidak ada. Di koordinat polar
$f=\cos2\theta/r^{2}$, dan pada seperempat cakram satuan
\[
\iint|f|\,dA\ge\int_{0}^{\pi/2}|\cos2\theta|\,d\theta
\int_{0}^{1}\frac{dr}{r}=\infty ,
\]
karena $\int_{0}^{1}dr/r$ menyimpang di $r=0$. Fubini tidak salah: syaratnya
tidak dipenuhi.

Kegagalan ketiga adalah batas yang tidak terurus. Gauss menuntut $E$ tertutup
dan terbatas dengan batas mulus sepotong-potong. Pada himpunan tak terbatas,
batasnya dapat lari ke tak hingga dan integral permukaannya berhenti berarti;
Spivak menaruh hal ini sebagai latihan setelah Teorema 5-5, yang gagal tanpa
kekompakan. Berhingga sudut masih aman, karena normalnya hanya melompat di
berhingga titik; batas bergerigi tanpa akhir, misalnya garis pantai, tidak.

Kegagalan keempat adalah permukaan yang tidak dapat diorientasikan. Pada pita
Mobius, berjalan mengelilingi pita membawa normal kembali ke sisi yang
berlawanan, sehingga $\mathbf{n}$ pada rumus Stokes tidak terdefinisi.
Teoremanya bukan salah; pernyataannya belum lengkap untuk benda semacam itu.

\begin{tantangan}
Ambil $T(u,v)=(u^{2},v)$ pada $D^{*}=(-2,2)\times(0,1)$. Peta ini tampak
melebarkan: titik $u=2$ mendarat di $x=4$, dan turunannya $2u$ tumbuh tanpa
batas ketika $u$ menjauh dari nol. Hitung luas $T(D^{*})$ dari gambarnya, lalu
hitung dengan rumus perubahan variabel. Kedua angka berbeda dengan faktor tepat
$2$. Jelaskan mengapa, dan sebutkan garis tempat determinan Jacobiannya nol.

Petunjuk 1: periksa apakah $T$ satu-satu. Ada dua nilai $u$ yang mendarat di
$x=1$.

Petunjuk 2: garis $u=0$ dilipat menjadi satu titik, sehingga tiap titik bayangan
terhitung dua kali oleh integral; $|\det J|=|2u|$ tidak tahu bahwa lipatan itu
hanya satu lembar.

Kalau mentok, kembali ke Definisi perubahan variabel yang sah dan periksa syarat
satu-satu.
\end{tantangan}

Satu peringatan penutup, yang tidak ada di buku teks. Kuadratur numerik tidak
pernah memeriksa syarat. Ia memberi $6.2831853072$ dengan tujuh digit benar
untuk medan yang melanggar syarat Green, dan tidak ada pesan galat. Bilangan yang rapi bukan bukti syaratnya dipenuhi; itu hanya bukti
aritmetikanya berjalan.

\section{Contoh hitung}
\tingkat{menengah}

Sekarang tiga hitungan yang menjadi tulang bab ini. Ketiganya bisa diperiksa
dengan Python, dan ketiganya pernah salah di tangan penulis sendiri.

\textbf{Satu integral, dua sistem koordinat.}
\nopagebreak[4]

Hitung
\[
I=\iint_{D}e^{-(x^{2}+y^{2})}\,dA,\qquad D=\{(x,y):x^{2}+y^{2}\le1\}.
\]
Di koordinat polar, $D$ menjadi $0\le r\le1$, $0\le\theta\le2\pi$, dan
$dA=r\,dr\,d\theta$. Integralnya terpisah menjadi dua faktor:
\[
I=\int_{0}^{2\pi}\!\!\int_{0}^{1}e^{-r^{2}}r\,dr\,d\theta
=2\pi\cdot\Bigl[-\frac12e^{-r^{2}}\Bigr]_{0}^{1}
=2\pi\cdot\frac{1-e^{-1}}{2}=\pi\bigl(1-e^{-1}\bigr).
\]
Nilainya $1.9858653038$. Antiturunan terhadap $r$ ada justru karena faktor
$r$: tanpa faktor itu, antiturunannya melibatkan fungsi galat.

Di koordinat Cartesius, daerahnya menuntut dua potongan:
\[
I=\int_{-1}^{1}\!\!\int_{-\sqrt{1-x^{2}}}^{\sqrt{1-x^{2}}}
e^{-(x^{2}+y^{2})}\,dy\,dx .
\]
Integral dalam tidak punya antiturunan elementer, jadi kita harus memakai
kuadratur. Hasil kedua bentuk, dengan toleransi yang sama:
\begin{align*}
\text{kartesius}&: 1.9858653037988705 \quad (11907\ \text{evaluasi}),\\
\text{polar}&: 1.9858653037988712 \quad (441).
\end{align*}
Galatnya $8.9\times10^{-16}$ dan $2.2\times10^{-16}$; selisih ketelitiannya
tidak penting, tetapi selisih biayanya $27$ kali.

Aturan titik tengah pada grid Cartesius yang sama rata menurunkan galat lebih
lambat: $1.141\times10^{-3}$ pada $50$ petak per sisi, dan
$4.667\times10^{-5}$ pada $400$. Setiap kali jumlah petak digandakan, galatnya
turun sekitar $2.9$ kali, bukan $4$ kali, karena tinggi kolom di dekat tepi
berperilaku seperti $\sqrt{1-x^{2}}$ sehingga galatnya seorde $h^{3/2}$.

\textbf{Volume bola dengan determinan Jacobian.}

Volume bola berjari-jari $R$ adalah
\[
V=\int_{0}^{2\pi}\!\!\int_{0}^{\pi}\!\!\int_{0}^{R}
\rho^{2}\sin\varphi\,d\rho\,d\varphi\,d\theta
=2\pi\cdot2\cdot\frac{R^{3}}{3}=\frac{4\pi R^{3}}{3}.
\]
Untuk $R=1$ nilainya $4.1887902048$. Kuadratur tiga dimensi pada bola memberi
\[
4.188790204786397 \quad(\text{selisih }6.2\times10^{-15}),
\]
dan simulasi Monte Carlo dengan $10^{6}$ titik pada kubus $[-1,1]^{3}$ memberi
$4.192824\pm0.004$.

Bandingkan dengan menghapus determinan Jacobian. Tanpa faktor
$\rho^{2}\sin\varphi$,
\[
\int_{0}^{2\pi}\!\!\int_{0}^{\pi}\!\!\int_{0}^{1}1\,d\rho\,d\varphi\,d\theta
=2\pi^{2}=19.7392088022,
\]
yaitu $3\pi/2$ kali terlalu besar, dan hasilnya tetap bilangan yang tampak
wajar untuk sebuah volume.



Tabel berikut mengumpulkan semua angka di bab ini.

\begin{center}
{\small
\begin{tabular}{llll}
\hline
Besaran & Eksak & Numerik & Galat \\
\hline
$\iint_{D}e^{-(x^{2}+y^{2})}dA$, polar & $1.985865$ & $1.985865$ & $2.2\times10^{-16}$ \\
$\iint_{D}e^{-(x^{2}+y^{2})}dA$, Cartesius & $1.985865$ & $1.985865$ & $8.9\times10^{-16}$ \\

$\int_{\gamma}F\cdot d\mathbf{r}$ & $10$ & $10.000000$ & $1.8\times10^{-15}$ \\
$\oint_{\partial D}$ Green, persegi & $2$ & $2.000000$ & $<10^{-15}$ \\
$\iint_{S}(\operatorname{curl}F)\cdot\mathbf{n}\,dA$ & $3.141593$ & $3.141593$ & $<10^{-13}$ \\
\hline
\end{tabular}}
\end{center}

\kertas{
Uji sendiri dalam tiga langkah. (1) Hitung ulang $I$ dengan tangan memakai
antiturunan $-\tfrac12e^{-r^{2}}$, lalu bandingkan dengan $1.9858653038$;
selisih di digit terakhir adalah galat pembulatan, bukan galat rumus. (2)
Jalankan \texttt{scipy.integrate.dblquad} pada kedua bentuk dan hitung berapa
kali fungsinya dipanggil; angka $11907$ dan $441$ akan muncul kembali, dengan
syarat toleransi keduanya sama. (3) Ganti $D$ menjadi cakram berjari-jari $2$
dan periksa bahwa jawaban eksaknya menjadi $\pi(1-e^{-4})$, yaitu $3.0841$.
Kalau langkah ketiga memberi angka lain, yang salah hampir pasti batas $r$,
bukan rumusnya. }

\section{Lab komputasi}
\tingkat{lanjut}

Kode berikut memeriksa lima hal sekaligus: biaya dua sistem koordinat, Green,
Gauss, Stokes, dan medan bersingularitas. Angka di tabel sebelumnya berasal
dari program ini.

\begin{lab}
\begin{verbatim}
import numpy as np
from scipy import integrate

# 1. Satu integral lipat dua, dua sistem koordinat
def fc(y, x):
    return np.exp(-(x*x + y*y))

cnt = [0]
def fc_count(y, x):
    cnt[0] += 1
    return fc(y, x)

kart = integrate.dblquad(fc_count, -1, 1,
                         lambda x: -np.sqrt(1-x*x),
                         lambda x: np.sqrt(1-x*x),
                         epsabs=1e-12, epsrel=1e-12)
nkart = cnt[0]
cnt[0] = 0

def fp(r, th):
    cnt[0] += 1
    return np.exp(-r*r)*r

pol = integrate.dblquad(fp, 0, 2*np.pi, 0, 1,
                        epsabs=1e-12, epsrel=1e-12)
nilai = np.pi*(1 - np.exp(-1))
print("kartesius %.16f galat %.2e eval %d"
      % (kart[0], abs(kart[0]-nilai), nkart))
print("polar     %.16f galat %.2e eval %d"
      % (pol[0], abs(pol[0]-nilai), cnt[0]))
print("eksak     %.16f" % nilai)

# 2. Green: (x^2-y)dx + (x+2y)dy pada persegi satuan
P = lambda x, y: x*x - y
Q = lambda x, y: x + 2*y
tepi = [integrate.quad(lambda t: P(t, 0.0), 0, 1)[0],
        integrate.quad(lambda t: Q(1.0, t), 0, 1)[0],
        integrate.quad(lambda t: P(t, 1.0), 1, 0)[0],
        integrate.quad(lambda t: Q(0.0, t), 1, 0)[0]]
ganda = integrate.dblquad(lambda y, x: 2.0, 0, 1, 0, 1)[0]
print("Green tepi %.15f  ganda %.15f" % (sum(tepi), ganda))

# 3. Gauss: F=(x^3,y^3,z^3) pada bola satuan
def fluks(phi, th):
    sp, cp = np.sin(phi), np.cos(phi)
    x, y, z = sp*np.cos(th), sp*np.sin(th), cp
    return (x**4 + y**4 + z**4)*sp

F = integrate.dblquad(fluks, 0, 2*np.pi, 0, np.pi,
                      epsabs=1e-11, epsrel=1e-11)[0]
print("Gauss fluks %.13f  12pi/5 %.13f" % (F, 12*np.pi/5))

# 4. Stokes: F=(z,x,y) pada belahan atas bola satuan
def curl_n(phi, th):
    sp, cp = np.sin(phi), np.cos(phi)
    x, y, z = sp*np.cos(th), sp*np.sin(th), cp
    return (x + y + z)*sp

S = integrate.dblquad(curl_n, 0, 2*np.pi, 0, np.pi/2)[0]
L = integrate.quad(lambda t: np.cos(t)**2, 0, 2*np.pi)[0]
print("Stokes permukaan %.13f  batas %.13f" % (S, L))

# 5. Medan dengan singularitas: F=(-y,x)/(x^2+y^2)
def lingkar(t):
    x, y = np.cos(t), np.sin(t)
    return (-y*(-np.sin(t)) + x*np.cos(t))/(x*x + y*y)

def Pg(x, y): return -y/(x*x + y*y)
def Qg(x, y): return x/(x*x + y*y)

h = 1e-5
curl_01 = ((Qg(1+h, 0) - Qg(1-h, 0))/(2*h)
           - (Pg(1, h) - Pg(1, -h))/(2*h))
C = integrate.quad(lingkar, 0, 2*np.pi)[0]
print("lingkaran %.13f  2pi %.13f" % (C, 2*np.pi))
print("curl(1,0) %.2e" % curl_01)
\end{verbatim}
\end{lab}

Keluaran sebenarnya, apa adanya:

\begin{verbatim}
kartesius 1.9858653037988705 galat 8.88e-16 eval 11907
polar     1.9858653037988712 galat 2.22e-16 eval 441
eksak     1.9858653037988714
Green tepi 2.000000000000000  ganda 2.000000000000000
Gauss fluks 7.5398223686155  12pi/5 7.5398223686155
Stokes permukaan 3.1415926535898  batas 3.1415926535898
lingkaran 6.2831853071796  2pi 6.2831853071796
curl(1,0) -2.01e-10
\end{verbatim}

Empat hal patut dibaca. Pertama, dua baris awal sama sampai digit ke-15 dengan
biaya berbeda $27$ kali: koordinat yang tepat tidak mengubah jawaban, hanya
harganya. Kedua, Green cocok sampai digit terakhir, tetapi itu
hanya bukti implementasinya konsisten dengan rumusnya. Ketiga, Gauss dan Stokes
dihitung dari sisi berbeda dan tetap bertemu; di situlah verifikasi punya gigi.
Keempat, kuadratur memberi $6.2831853071796$ untuk medan yang melanggar syarat
Green, tanpa satu pun pesan galat.

\begin{untukdosen}
Urutan mengajar yang saya pakai: kerapatan pelat dan program yang lupa faktor
$r$, lalu Fubini, lalu determinan Jacobian sebagai luas petak yang meregang,
lalu integral garis dan teorema dasarnya, baru Green. Green sebaiknya dibuktikan
penuh untuk daerah bertipe I di papan, karena bukti itu hanya memakai dua alat
yang sudah mereka punya. Gauss dan Stokes cukup dibuktikan untuk kotak dan
permukaan berparameter.

Dua miskonsepsi yang paling sering muncul. Pertama, mahasiswa menganggap
$\partial P/\partial y=\partial Q/\partial x$ cukup untuk konservatif, karena
mereka hanya menguji medan pada domain tanpa lubang. Medan
$(-y,x)/(x^{2}+y^{2})$ harus muncul di papan, bukan di daftar bacaan. Kedua,
mereka memperlakukan determinan Jacobian sebagai formalitas notasi dan menulis
$r$ atau $\rho^{2}\sin\varphi$ sebagai hafalan tanpa memeriksa batas.
Bandingkan $4\pi/3$ dengan $2\pi^{2}$; rasio $4.712$ kali itu lebih mengesankan
daripada seratus kali peringatan.

Satu aktivitas kelas: empat kelompok, masing-masing menghitung satu ruas dari
Green, Gauss, dan Stokes untuk medan yang sama, lalu mencocokkan angka. Kelompok
yang salah tanda ketahuan dalam lima menit, dan itu lebih baik daripada diberi
tahu. Tanda kesiapan lanjut: mahasiswa mampu menjelaskan mengapa orientasi
mengubah tanda tanpa melihat rumus.
\end{untukdosen}

\section{Ringkasan}
\begin{ringkasan}
\begin{itemize}
\item Integral lipat dua didefinisikan lewat jumlah atas dan bawah pada petak;
Fubini mengizinkan urutan pengintegralan ditukar bila fungsinya terintegralkan.
\item Mengganti koordinat mengalikan luas petak dengan $|\det J|$; polar memberi
$r$, bola memberi $\rho^{2}\sin\varphi$. Tanpa faktor itu hasilnya rapi tetapi
salah.
\item Pada medan konservatif, integral garis hanya bergantung pada titik ujung.
\item Green, Gauss, dan Stokes menyatakan satu hal: integral turunan pada daerah
sama dengan integral medan pada batasnya, dan orientasi menentukan tanda.
\item Semua teorema itu bersyarat: singularitas, batas tak terurus, permukaan tak
terorientasi, dan fungsi tak terintegralkan membuatnya gagal tanpa peringatan.
\end{itemize}
\end{ringkasan}

\section{Latihan}
\begin{latihan}
\begin{enumerate}
\item Hitung $\iint_{D}(x^{2}+y^{2})\,dA$ pada cakram satuan dengan Cartesius dan polar, periksa dengan \texttt{dblquad}, dan sebutkan jumlah evaluasi tiap cara.

\item Buktikan bahwa bila $T(u,v)=A(u,v)$ linear dan $\det A\ne0$, rumus
perubahan variabel menjadi $|\det A|\iint_{D}f\circ T\,dA$.

\item Ulangi contoh Green pada persegi dengan arah jarum jam: tunjukkan hasilnya
$-2$ dan sebutkan sisi mana menyumbang tanda mana.

\item Hitung volume elipsoid $x^{2}/a^{2}+y^{2}/b^{2}+z^{2}/c^{2}\le1$ dengan
perubahan variabel $x=au$, $y=bv$, $z=cw$, lalu periksa rumusnya pada
$a=b=c=1$.

\item Nilai klaim ini: benar, salah, atau belum cukup bukti. ``Karena
$\operatorname{curl}F=0$ di setiap titik tempat $F$ terdefinisi, medan
$F=(-y,x)/(x^{2}+y^{2})$ konservatif pada
$\mathbb{R}^{2}\setminus\{(0,0)\}$.'' Sertakan alasan dan satu angka
pendukung.

\item Buktikan bahwa integral garis tidak bergantung lintasan jika dan hanya jika
sirkulasi setiap kurva tertutup nol, dengan memakai dua lintasan dari $A$ ke
$B$ yang membentuk satu kurva tertutup.

\item Hitung kedua ruas Stokes untuk $F=(-y,x,z)$ pada paraboloid
$z=1-x^{2}-y^{2}$, $z\ge0$, dengan normal mengarah ke atas, lalu periksa
dengan kuadratur.

\item Hitung kedua ruas Gauss untuk $F=(x,y,z)/(x^{2}+y^{2}+z^{2})^{3/2}$ pada bola
satuan. Jelaskan mengapa hasilnya tidak bertentangan dengan teorema, dan
syarat mana yang dilanggar. (Petunjuk: hitung $\operatorname{div}F$ lebih dulu.)

\item Tunjukkan bahwa determinan Jacobian koordinat bola nol pada $\varphi=0$ dan
$\varphi=\pi$, lalu jelaskan mengapa peta itu tetap boleh dipakai pada
$0<\varphi<\pi$. (Petunjuk: periksa di mana peta berhenti satu-satu.)
\end{enumerate}
\end{latihan}

\begin{untukpeneliti}
Yang benar-benar terbukti di bab ini, dengan bukti yang ditulis: lemma daerah
di antara dua grafik, teorema dasar untuk integral garis, kebalikannya lewat
potensial integral, Green untuk daerah bertipe I dan II, akibat uji curl pada
domain terhubung sederhana, dan sketsa bukti Gauss lewat kotak serta Stokes
lewat parametrisasi. Yang saya kutip tanpa bukti: Fubini (Spivak, bab
\emph{Integration}, Teorema 3-10), perubahan variabel penuh (Spivak bab yang
sama; Rudin, bab \emph{Integration of Differential Forms}), Gauss (Spivak,
Teorema 5-8), dan Stokes (Spivak, Teorema 5-9); Stokes umum adalah Teorema
4-13 dan 5-5.

Yang dipopulerkan secara berlebihan adalah ungkapan ``curl nol berarti
konservatif''. Yang benar: curl nol berarti bentuk tertutup, dan bentuk tertutup
eksak hanya bila grup kohomologi pertama domainnya nol. Untuk $\mathbb{R}^{3}$
itu berlaku, untuk $\mathbb{R}^{2}\setminus\{0\}$ tidak, dan selisihnya
diukur oleh integral pada lingkaran, yaitu $2\pi$ di sini.

Pertanyaan terbuka yang sungguh terbuka bagi saya: dapatkah kuadratur memberi
sertifikat bahwa syarat teoremanya dipenuhi, bukan hanya taksiran nilainya?
Keluaran di lab menunjukkan bahwa jawaban numerik yang benar sampai lima belas
digit dapat berasal dari medan yang melanggar syarat, tanpa tanda bahaya. Untuk
medan yang hanya diketahui lewat sampel pada grid, tidak ada cara umum
memastikan tidak ada singularitas di antara titik sampel.

Batas keberlakuan: semua fungsi di sini bernilai real dan mulus
sepotong-sepotong; semua angka dari Python 3.14.4, NumPy 2.5.2, SciPy 1.18.0
dengan toleransi $10^{-12}$ kecuali disebut lain, dan permukaan bola
diparametrisasi dengan $\mathbf{n}$ mengarah keluar.
\end{untukpeneliti}

\peta{Bab 15 pada Stewart, \emph{Calculus: Early Transcendentals}, bab ``Multiple Integrals'' (15.1--15.3 dan 15.6--15.9) dan bab ``Vector Calculus'' (16.2
integral garis, 16.3 teorema dasar, 16.4 Green, 16.5 curl
dan divergensi, 16.8 Stokes, 16.9 Gauss); pada Spivak, \emph{Calculus on
Manifolds}, bab ``Integration'' (Fubini, Teorema 3-10, perubahan variabel),
bab ``Integration on Chains'' dan bab ``Integration of Manifolds'' (Teorema
5-5, 5-8, 5-9); pada Rudin, \emph{Principles of Mathematical Analysis}, bab
``Functions of Several Variables'' dan ``Integration of Differential Forms''; pada Active Calculus, bab 11 dan 12 untuk urutan yang lebih bertahap.
Basisnya bab \emph{Integral, Deret Fungsi, dan Pertukaran Limit} (bagian integral)
dan bab \emph{Kekontinuan dan Turunan} (turunan parsial, Jacobian) pada edisi pertama
buku ini. Bab ini memakai integral Riemann dan teorema dasar kalkulus
dari Bab 13, kekontinuan seragam dan himpunan kompak dari Bab 11, determinan
dari Bab 7, serta gradien dan matriks Jacobian dari Bab 12, semuanya dirujuk
dan tidak diulang.}
